\documentclass[11pt]{article}

\usepackage{amsmath}
\usepackage{amssymb}
\usepackage[margin=1in]{geometry}
\usepackage{hyperref}

\title{Kernel Modeling of Nuisance Ratios for Systematics-Aware NPLM}
\author{}
\date{}

\begin{document}
\maketitle

\section{Goal}

The baseline kernel NPLM learns a flexible log-density ratio between data and a
central reference sample. With nuisance parameters, the null hypothesis is
composite. We write the nuisance-deformed reference density as
\begin{equation}
  n(x \mid R_\nu) = r(x;\nu)\,n(x \mid R_0),
\end{equation}
where $R_0$ is the central reference. The full alternative model can then be
written as
\begin{equation}
  h(x;\alpha,\nu) = f_K(x;\alpha) + \log r(x;\nu),
\end{equation}
where $f_K$ is the residual kernel NPLM component and $\log r(x;\nu)$ is the
systematic deformation.

The first implementation step is to model $\log r(x;\nu)$ with kernel methods
without modifying the existing baseline NPLM implementation.

\section{Option A: Finite-Difference Falkon Morphing}

Assume that nuisance-varied Monte Carlo samples are available at
$\nu=+\epsilon$ and $\nu=-\epsilon$, together with a central sample at $\nu=0$.
Train two LogisticFalkon density-ratio estimators:
\begin{align}
  g_+(x) &\simeq \log \frac{p(x\mid +\epsilon)}{p(x\mid 0)}, \\
  g_-(x) &\simeq \log \frac{p(x\mid -\epsilon)}{p(x\mid 0)}.
\end{align}
The local nuisance response is then estimated by the centered finite
difference
\begin{equation}
  \delta(x) =
  \frac{g_+(x)-g_-(x)}{2\epsilon}.
\end{equation}
For one nuisance parameter, the morphing model is
\begin{equation}
  \log r(x;\nu) \simeq \nu\,\delta(x).
\end{equation}
For several independent nuisance directions, the first-order extension is
\begin{equation}
  \log r(x;\nu) \simeq \sum_a \nu_a\,\delta_a(x).
\end{equation}

This is the closest kernel analogue of the neural-network nuisance setup:
the nuisance log-ratio is a separate component, while the final NPLM fit learns
only the residual discrepancy $f_K(x)$.

\subsection{Classifier Weighting}

For central samples $x\sim p_0$ and varied samples $x\sim p_\nu$, a weighted
binary logistic classifier has population minimizer
\begin{equation}
  f^*(x)=
  \log \frac{N_\nu p_\nu(x)}
            {w_0 N_0 p_0(x)}.
\end{equation}
To learn the normalized shape ratio $p_\nu/p_0$, choose
\begin{equation}
  w_0 = \frac{N_\nu}{N_0}.
\end{equation}
If the nuisance also changes the expected total rate by a factor
$\rho_\nu=N_\nu^{\mathrm{expected}}/N_0^{\mathrm{expected}}$, choose
\begin{equation}
  w_0 = \frac{N_\nu}{N_0\,\rho_\nu},
\end{equation}
so that the score estimates
\begin{equation}
  \log \rho_\nu \frac{p_\nu(x)}{p_0(x)}.
\end{equation}

\subsection{Caching}

Falkon does not currently provide a simple native model serialization path.
For the profiled NPLM statistic this is not a blocker, because the likelihood
only needs repeated evaluations of $\log r(x;\nu)$ on fixed event arrays.
After training the ratio estimators, cache
\begin{align}
  \delta_{\mathrm{ref},i} &= \delta(x^{\mathrm{ref}}_i), \\
  \delta_{\mathrm{data},j} &= \delta(x^{\mathrm{data}}_j),
\end{align}
and save these arrays as a NumPy archive. Profiling then evaluates
\begin{align}
  \log r(x^{\mathrm{ref}}_i;\nu) &= \nu\,\delta_{\mathrm{ref},i}, \\
  \log r(x^{\mathrm{data}}_j;\nu) &= \nu\,\delta_{\mathrm{data},j}.
\end{align}
This makes the nuisance layer cheap and independent of Falkon model persistence.

\section{One-Dimensional Validation}

The first validation toy uses a truncated exponential density on $[0,x_{\max}]$,
\begin{equation}
  p(x\mid \lambda) =
  \frac{\lambda e^{-\lambda x}}
       {1-e^{-\lambda x_{\max}}}.
\end{equation}
The nuisance changes the rate as
\begin{equation}
  \lambda(\nu)=\lambda_0 e^{a\nu}.
\end{equation}
The analytic response at $\nu=0$ is
\begin{equation}
  \delta_{\mathrm{true}}(x)
  =
  a - a\lambda_0 x
  - a\lambda_0
    \frac{x_{\max} e^{-\lambda_0 x_{\max}}}
         {1-e^{-\lambda_0 x_{\max}}}.
\end{equation}
The implementation compares the Falkon finite-difference estimate
$\delta(x)$ against this analytic curve.

\section{Option B: Conditional Kernel Ratio}

An alternative is to train a single conditional model on augmented inputs
\begin{equation}
  z=(x,\nu)
\end{equation}
so that
\begin{equation}
  g(x,\nu) \simeq \log \frac{p(x\mid \nu)}{p(x\mid 0)}.
\end{equation}
This can represent nonlinear nuisance dependence without a Taylor expansion.
However, it requires careful kernel choices in the joint $(x,\nu)$ space,
adequate coverage of the nuisance grid, and validation of interpolation in
$\nu$. For this reason the first implementation uses option A.

\section{Current Implementation Boundary}

The initial code lives in a separate \texttt{nplm\_systematics} package. It
does not modify the baseline \texttt{nplm.LogFalkonNPLM} class. The current
scope is only the nuisance-ratio model
\begin{equation}
  \log r(x;\nu),
\end{equation}
not yet the full profiled statistic $t=\tau-\Delta$.

\end{document}
